
\documentclass[11pt]{article}

\usepackage[a4paper,margin=2.05cm]{geometry}
\usepackage[T1]{fontenc}
\usepackage[utf8]{inputenc}
\usepackage[english]{babel}
\usepackage{amsmath,amssymb,amsthm,mathtools,bm}
\usepackage{booktabs,longtable,array}
\usepackage{graphicx}
\usepackage{microtype}
\usepackage{enumitem}
\usepackage{xcolor}
\usepackage[hidelinks]{hyperref}
\usepackage{caption}
\usepackage{float}

\newtheorem{theorem}{Theorem}[section]
\newtheorem{proposition}[theorem]{Proposition}
\newtheorem{lemma}[theorem]{Lemma}
\newtheorem{corollary}[theorem]{Corollary}
\theoremstyle{definition}
\newtheorem{definition}[theorem]{Definition}
\newtheorem{remark}[theorem]{Remark}
\newtheorem{problem}[theorem]{Problem}

\newcommand{\dd}{\,\mathrm d}
\newcommand{\R}{\mathbb R}
\newcommand{\C}{\mathbb C}
\newcommand{\Ree}{\operatorname{Re}}
\newcommand{\Imm}{\operatorname{Im}}
\newcommand{\sgn}{\operatorname{sgn}}

\title{\textbf{Euclidean Lobe Geometry, Reciprocal Scattering, and Transverse Dynamics of the Riemann Zeta Function}}
\author{Víctor Casaña-Giner\\
\small Bio Logic Crop Science S.L.\\
\small ORCID: 0000-0002-8252-2608}
\date{24 August 2026}

\begin{document}
\maketitle

\begin{abstract}
We treat the Euclidean trace
\[
\Gamma(t)=\left(t,\Re\zeta\!\left(\tfrac12+it\right),
                 \Im\zeta\!\left(\tfrac12+it\right)\right)
\]
as a primary geometric object.  The trace, referred to throughout as the
\emph{Violet Thread}, is naturally segmented by its successive contacts with
the value-axis origin into zero-to-zero lobes.  The construction packages,
within the same object, radial excursion, angular sweep, Euclidean length,
zero-to-zero segmentation, and an exact transverse quarter-turn law.

The central structural equation is global and first order.  With transverse
coordinate $\delta$,
\[
F(\delta,t)=\zeta\!\left(\tfrac12+\delta+it\right)
\]
satisfies
\[
F_\delta=-iF_t.
\]
Thus the entire transverse family is generated from the longitudinal geometry
by an exact $90^\circ$ rotation of the value-plane velocity.  This is the
complete first-order evolution law, not the first term of an independent
higher-order hierarchy.

The visually striking feature of the thread is that its zero-to-zero lobes
are eccentric: the zero-axis is a contact/pinching locus, not the geometric
center of the lobes.  We show that the smooth Riemann--von Mangoldt/Lambert
phase produces the corresponding ideal geometry exactly.  A unit phase sweep
with a one-arch radial profile is an eccentric circle tangent to the origin.
The exact RVM/Lambert model consists of such tangent circles and satisfies a
global first-order no-exit theorem with an explicit positive modulus barrier
for every nonzero transverse displacement.

For the actual zeta lobes we introduce a phase centroid $\mathcal C_n$,
normalized eccentricity $E_n$, and circle-tangency defect $D_n$.  On the first
49 zero intervals, the measured $E_n$ is predicted by the closed-form
RVM-phase one-arch law with correlation $0.997589$, RMSE $0.017527$, and MAE
$0.014155$.  The same lobe architecture also controls Euclidean arc inflation:
directly recomputed $R_n=L_n/(\gamma_{n+1}-\gamma_n)$ agrees with its
one-arch elliptic reduction at correlation $0.999162$.  Global phase/RVM
diagnostics use all $2{,}001{,}052$ supplied zeta and Lambert ordinates.

The paper isolates the analytic issue cleanly.  A global no-exit theorem for
the exact zeta lobes, derived from the same first-order field and their
eccentric geometry, would already give zero confinement.  No subsequent
second-, third-, or infinite-order condition would remain.
\end{abstract}

\tableofcontents
\newpage

\section{Main geometric statements}

The purpose of this section is to state the architecture before introducing
technical details.

\begin{proposition}[The Violet Thread packages five geometries]\label{prop:package}
Let
\[
W(t)=\zeta\!\left(\frac12+it\right).
\]
The single Euclidean curve
\[
\Gamma(t)=(t,\Re W(t),\Im W(t))
\]
simultaneously carries:
\begin{enumerate}[leftmargin=1.8em]
\item the radial excursion $|W(t)|$;
\item the angular sweep $\arg W(t)$ between successive zeros;
\item the Euclidean arc element
\[
\dd\ell=\sqrt{1+|W_t(t)|^2}\,\dd t;
\]
\item the zero-to-zero decomposition
\[
[\gamma_n,\gamma_{n+1}];
\]
\item the generator of the complete transverse continuation,
\[
F_\delta=-iF_t.
\]
\end{enumerate}
Thus the longitudinal lobe geometry and the transverse generator are not two
separate constructions.
\end{proposition}

\begin{theorem}[Exact global quarter-turn law]\label{thm:quarter-main}
For
\[
F(\delta,t)=\zeta\!\left(\frac12+\delta+it\right),
\]
one has everywhere in the analytic strip
\[
\boxed{F_\delta=-iF_t.}
\]
Consequently the transverse and longitudinal value-plane velocities are
orthogonal and have identical magnitude.
\end{theorem}

\begin{theorem}[First-order sufficiency of global confinement]\label{thm:firstorder-suff}
Suppose the exact zeta flow governed by Theorem~\ref{thm:quarter-main} is
shown, by a global property of this first-order field and the exact lobe
geometry, to satisfy
\[
F(\delta,t)\neq0
\qquad
\text{for every }t\in\R,\quad 0<|\delta|<\frac12.
\]
Then every nontrivial zero is confined to the central thread.  No independent
second-, third-, or higher-order transverse condition is required.
\end{theorem}

\begin{proof}
The displayed nonvanishing statement is exactly zero confinement in the
transverse coordinate.  Equation $F_\delta=-iF_t$ is already the complete
first-order evolution equation for $F$; higher $\delta$-derivatives are
obtained by differentiating it and are not additional hypotheses.
\end{proof}

\begin{theorem}[RVM tangent-circle no-exit]\label{thm:rvm-main}
The smooth RVM/Lambert tangent-circle model defined in
Section~\ref{sec:rvm} satisfies the same first-order translation law and,
for $t>2\pi$, cannot meet the origin at any nonzero transverse displacement.
An explicit positive lower bound for its modulus is given in
Theorem~\ref{thm:rvmnoexit}.
\end{theorem}

These statements separate what is exact from what remains to be transferred.
The paper does not replace the Violet Thread by an external zero criterion:
the unresolved analytic content is whether the exact zeta lobes inherit a
global no-exit property from the eccentric first-order geometry quantified
below.

\section{Notation}

\begin{center}
\begin{tabular}{>{\raggedright\arraybackslash}p{0.22\textwidth}p{0.70\textwidth}}
\toprule
Symbol & Meaning \\
\midrule
$W(t)$ & $\zeta(\tfrac12+it)$, complex value trace of the Violet Thread \\
$\Gamma(t)$ & $(t,\Re W(t),\Im W(t))\in\R^3$ \\
$Z(t)$ & Hardy's real function, $W(t)=Z(t)e^{-i\vartheta(t)}$ \\
$\vartheta(t)$ & Riemann--Siegel phase \\
$\gamma_n$ & successive zero ordinates of $W$ \\
$I_n$ & zero-to-zero interval $(\gamma_n,\gamma_{n+1})$ \\
$g_n$ & $\gamma_{n+1}-\gamma_n$ \\
$\Omega_n$ & $\vartheta(\gamma_{n+1})-\vartheta(\gamma_n)$ \\
$s_n$ & $\Omega_n/\pi$, phase-unfolded lobe sweep \\
$A_n$ & $\max_{I_n}|Z(t)|$ \\
$L_n$ & Euclidean arc length of the $n$th lobe \\
$R_n$ & $L_n/g_n$ \\
$F(\delta,t)$ & $\zeta(\tfrac12+\delta+it)$, transverse continuation \\
$\Theta_0$ & $\frac z2\log(z/(2\pi e))-\pi/8$, smooth RVM phase \\
$\tau_n$ & supplied RVM/Lambert ordinates \\
$\mathcal C_n$ & phase centroid of a rotated, sign-normalized lobe \\
$E_n$ & normalized eccentricity $2|\mathcal C_n|/A_n$ \\
$D_n$ & normalized circle-tangency defect \\
\bottomrule
\end{tabular}
\end{center}

\section{The Euclidean thread itself}

\begin{definition}[Violet Thread]
Define
\[
\boxed{
W(t):=\zeta\!\left(\frac12+it\right)=u(t)+iv(t)
}
\]
and
\[
\boxed{
\Gamma(t)=\bigl(t,u(t),v(t)\bigr).
}
\]
\end{definition}

The height $t$ is retained as an actual Euclidean coordinate.  A crossing in
the value-plane projection is therefore not a self-intersection of
$\Gamma\subset\R^3$ unless the heights also agree.

\begin{figure}[H]
\centering
\includegraphics[width=.97\textwidth]{../figures/fig01_euclidean_zeta_thread.png}
\caption{The Euclidean zeta thread over the first fifty supplied zero
ordinates.  The lobes repeatedly contact the value-axis origin while their
bulk lies away from that axis.}
\end{figure}

\begin{figure}[H]
\centering
\includegraphics[width=.80\textwidth]{../figures/fig02_transverse_eccentric_lobes.png}
\caption{Transverse projection.  The common zero point is not the center of
the lobe geometry; it is an eccentric contact/pinching point.}
\end{figure}

This eccentricity is not an incidental plotting effect.  It is forced by the
phase-amplitude structure described next.

\section{Exact cylindrical anatomy of a zero-to-zero lobe}

Hardy's function gives
\[
\boxed{
W(t)=Z(t)e^{-i\vartheta(t)}.
}
\]
Hence
\[
\boxed{
\Gamma(t)=
\bigl(t,Z(t)\cos\vartheta(t),-Z(t)\sin\vartheta(t)\bigr).
}
\]
Differentiating,
\[
W_t=e^{-i\vartheta}(Z'-i\vartheta'Z),
\]
and therefore
\[
\boxed{
|W_t|^2=Z'^2+\vartheta'^2Z^2.
}
\]

The decomposition is geometrically exact:
\[
\boxed{
\text{radial speed}=Z',
\qquad
\text{angular speed contribution}=\vartheta'Z.
}
\]

Between two successive zeros the sign of $Z$ is constant.  Put
\[
\varepsilon_n=\sgn Z(t),\qquad t\in I_n,
\]
and rotate away the absolute initial phase:
\[
\widehat W_n(t)
=
\varepsilon_ne^{i\vartheta(\gamma_n)}W(t).
\]
Then
\[
\boxed{
\widehat W_n(t)
=
|Z(t)|e^{-i\phi_n(t)},
\qquad
\phi_n(t)=\vartheta(t)-\vartheta(\gamma_n).
}
\]
Thus every lobe has a nonnegative radial profile and a phase sweep.  Define
\[
\boxed{
\Omega_n=\phi_n(\gamma_{n+1}),
\qquad
s_n=\frac{\Omega_n}{\pi}.
}
\]

This $s_n$ is not an auxiliary statistic attached after the fact: it is the
angular extent of the actual Euclidean lobe.

\section{Why the zeros are eccentric}

Introduce uniform phase coordinate
\[
q=\frac{\phi_n(t)}{\Omega_n},
\qquad 0\le q\le1.
\]
The simplest single-arch radial profile is
\[
r_n(q)=A_n\sin(\pi q).
\]
Keeping the actual lobe phase sweep gives
\[
\boxed{
\widehat W_n^{(\mathrm{arch})}(q)
=
A_n\sin(\pi q)e^{-i\pi s_nq}.
}
\]

\begin{theorem}[Unit sweep produces a tangent circle]\label{thm:circle}
When $s_n=1$,
\[
\widehat W_n^{(\mathrm{arch})}(q)
=
-\frac{iA_n}{2}
+\frac{iA_n}{2}e^{-2\pi iq},
\]
so
\[
\boxed{
\left|
\widehat W_n^{(\mathrm{arch})}(q)+\frac{iA_n}{2}
\right|
=
\frac{A_n}{2}.
}
\]
The lobe is exactly a circle whose center is displaced by one radius from the
origin.  The origin is a tangency point.
\end{theorem}

\begin{proof}
Use
\[
\sin(\pi q)e^{-i\pi q}
=
\frac{1-e^{-2\pi iq}}{2i}.
\]
\end{proof}

This elementary identity explains the systematic visual eccentricity:
a one-arch radial pulse combined with one RVM phase unit cannot be centered
on the zero axis.  Its natural center is displaced, and the zero is the point
where the displaced circle touches the origin.

\begin{figure}[H]
\centering
\includegraphics[width=.72\textwidth]{../figures/fig20_tangent_circle_identity.png}
\caption{The unit-sweep one-arch identity.  Rotation changes orientation but
not the fact that the circle is eccentric and tangent to the origin.}
\end{figure}

The exact zeta lobe is not assumed to be a perfect circle.  The point is that
the dominant phase/radial architecture already predicts why the zeros occupy
an eccentric contact locus.

\begin{figure}[H]
\centering
\includegraphics[width=.75\textwidth]{../figures/fig17_single_lobe_phase_geometry.png}
\caption{Representative actual lobe in uniform phase coordinate together
with its one-arch phase model and phase centroid.}
\end{figure}

\section{The RVM/Lambert phase creates the ideal lobe geometry}\label{sec:rvm}

The smooth counting phase is
\[
\boxed{
\Theta_0(z)
=
\frac z2\log\frac{z}{2\pi e}-\frac{\pi}{8}.
}
\]
The supplied Lambert ordinates satisfy
\[
\boxed{
\Theta_0(\tau_n)
=
\pi\left(n-\frac32\right),
}
\]
so each consecutive interval carries exactly one phase unit $\pi$.

Define
\[
\boxed{
M_0(t)=\cos\Theta_0(t)e^{-i\Theta_0(t)}.
}
\]
Then
\[
M_0(t)
=
\frac12\left(1+e^{-2i\Theta_0(t)}\right),
\]
and therefore
\[
\boxed{
\left|M_0(t)-\frac12\right|=\frac12.
}
\]
Every RVM/Lambert zero-to-zero lobe is an exact circle tangent to the origin.

This gives a precise interpretation:

\[
\boxed{
\text{RVM supplies the angular tangent-circle skeleton;}
\quad
\text{the actual }Z\text{-profile supplies the arithmetic deformation.}
}
\]

\begin{figure}[H]
\centering
\includegraphics[width=.94\textwidth]{../figures/fig13_rvm_tangent_circle_thread.png}
\caption{Exact tangent-circle thread generated by the RVM/Lambert phase.}
\end{figure}

\begin{figure}[H]
\centering
\includegraphics[width=.98\textwidth]{../figures/fig16_zeta_rvm_overlay_3d.png}
\caption{Actual zeta lobes and the RVM/Lambert tangent-circle skeleton shown
in the same Euclidean coordinates.  The comparison is geometric, not merely
a comparison of zero ordinates.}
\end{figure}

\section{Exact first-order transverse dynamics}

Set
\[
F(\delta,t)=
\zeta\!\left(\frac12+\delta+it\right),
\qquad
F(0,t)=W(t).
\]

\begin{proof}[Proof of Theorem~\ref{thm:quarter-main}]
Writing $s=\frac12+\delta+it$,
\[
F_\delta=\zeta'(s),
\qquad
F_t=i\zeta'(s),
\]
hence
\[
F_\delta=-iF_t.
\]
\end{proof}

The crucial point is that this is not a local approximation at $\delta=0$.
It is the exact PDE obeyed by the complete two-variable family.

On the central thread,
\[
F_\delta(0,t)
=
-iW_t(t)
=
e^{-i\vartheta(t)}
\left(-\vartheta'(t)Z(t)-iZ'(t)\right).
\]
Hence, in the Hardy frame,
\[
\boxed{
N_{\rm radial}=-\vartheta'Z,
\qquad
N_{\rm angular}=-Z'.
}
\]

The same lobe therefore carries its transverse generator.  Radial position,
radial derivative and RVM angular velocity determine how the lobe initially
moves across the transverse coordinate.

\begin{figure}[H]
\centering
\includegraphics[width=.75\textwidth]{../figures/fig18_quarter_turn_vector_field.png}
\caption{Along an actual lobe, each longitudinal value-plane tangent has an
exact quarter-turned transverse direction.  These vectors illustrate the
global law $F_\delta=-iF_t$ on the central trace.}
\end{figure}

\subsection{Why no hierarchy of higher-order conditions is needed}

Differentiating the first-order PDE gives identities such as
\[
F_{\delta\delta}=-F_{tt},
\qquad
F_{\delta\delta\delta}=iF_{ttt},
\]
but these are consequences of the same evolution equation.  They do not
constitute separate dynamical laws.

Equivalently, if
\[
\mathcal W(z)=\zeta\!\left(\frac12+iz\right),
\]
then
\[
\boxed{
F(\delta,t)=\mathcal W(t-i\delta).
}
\]
The transverse family is the analytic continuation of the same thread data.
There is no independent ``off-thread'' degree of freedom to be specified.

\begin{proposition}[Single-thread determination]\label{prop:identity}
If another analytic family agrees with $W(t)$ on an interval of the central
thread and is analytic in the same connected domain, then it agrees with the
corresponding continuation everywhere in that domain.
\end{proposition}

\begin{proof}
This is the identity theorem applied to the analytic continuation in the
complexified thread coordinate.
\end{proof}

This proposition is structural: it says that the transverse geometry is
encoded in the thread.  It does not by itself prove nonvanishing; the no-exit
content must come from a global geometric property of that encoded flow.

\section{Conformal and minimal surface generated by the thread}

Write $F=u+iv$ and define
\[
X(\delta,t)=(\delta,t,u(\delta,t),v(\delta,t))\in\R^4.
\]
Cauchy--Riemann gives
\[
u_\delta=v_t,\qquad v_\delta=-u_t.
\]
Consequently
\[
\boxed{X_\delta\cdot X_t=0}
\]
and
\[
\boxed{
\|X_\delta\|^2=\|X_t\|^2
=
1+\left|\zeta'\!\left(\frac12+\delta+it\right)\right|^2.
}
\]
Thus the induced metric is
\[
\boxed{
\dd s^2
=
\left(1+|\zeta'|^2\right)
(\dd\delta^2+\dd t^2).
}
\]
The holomorphic graph is conformal and minimal.  Its Gaussian curvature is
\[
\boxed{
K
=
-\frac{2|\zeta''|^2}{(1+|\zeta'|^2)^3}
\le0.
}
\]

The central zero-to-zero arc length introduced below is therefore not an
arbitrary statistic: it is the intrinsic metric length of a coordinate curve
on the same minimal surface whose transverse direction is generated by the
quarter-turn law.

\subsection{Intersection orientation}

Let the value-axis plane in the four-dimensional graph be
\[
\mathcal A=\{(\delta,t,0,0)\}.
\]
Zeros are intersections of the zeta surface with $\mathcal A$.  The
transverse Jacobian is
\[
\det
\begin{pmatrix}
u_\delta&u_t\\
v_\delta&v_t
\end{pmatrix}
=
u_t^2+v_t^2
=
|F_t|^2.
\]
At a simple zero this is strictly positive.  Thus every simple intersection
has the same local orientation.  This is a geometric property of the same
surface, not a separate zero criterion.

\section{Polar form of the transverse evolution}

Away from zeros write
\[
F(\delta,t)=r(\delta,t)e^{i\alpha(\delta,t)}.
\]
The first-order quarter-turn law is equivalent to
\[
\boxed{
(\log r)_\delta=\alpha_t,
\qquad
\alpha_\delta=-(\log r)_t.
}
\]
On the central thread, between zeros,
\[
\alpha(0,t)=-\vartheta(t)+k\pi,
\]
and therefore
\[
\boxed{
(\log r)_\delta(0,t)=-\vartheta'(t),
}
\]
while
\[
\boxed{
\alpha_\delta(0,t)=-\frac{Z'(t)}{Z(t)}.
}
\]

These identities provide a direct geometric conjugacy:
the longitudinal twist of the lobe controls normal radial response, while the
longitudinal radial profile controls normal phase response.

\section{Exact no-exit in the RVM tangent-circle model}

Analytically continue the RVM skeleton by
\[
M(\delta,t)=M_0(t-i\delta)
=
\frac12\left(
1+e^{-2i\Theta_0(t-i\delta)}
\right).
\]
It satisfies
\[
M_\delta=-iM_t.
\]

\begin{theorem}[Global RVM no-exit and modulus barrier]\label{thm:rvmnoexit}
For $t>2\pi$ and $\delta\ne0$,
\[
\boxed{M(\delta,t)\ne0.}
\]
Moreover
\[
\boxed{
|M(\delta,t)|
\ge
\frac12
\left|
1-\left(\frac{t}{2\pi}\right)^{-\delta}
\right|
>0.
}
\]
\end{theorem}

\begin{proof}
Let
\[
v(\delta)=\Im\Theta_0(t-i\delta).
\]
Since
\[
\Theta_0'(z)=\frac12\log\frac{z}{2\pi},
\]
\[
v'(\delta)
=
-\frac12\log
\frac{\sqrt{t^2+\delta^2}}{2\pi}.
\]
For $t>2\pi$ this is strictly negative for every real $\delta$, and $v(0)=0$.
Hence $v(\delta)$ has the opposite sign to $\delta$ and is nonzero whenever
$\delta\ne0$.

If $M(\delta,t)=0$, then
\[
e^{-2i\Theta_0(t-i\delta)}=-1,
\]
whose modulus is one.  But
\[
\left|e^{-2i\Theta_0(t-i\delta)}\right|
=
e^{2v(\delta)}\ne1
\]
for nonzero $\delta$, contradiction.

For the lower bound use
\[
|M|
\ge
\frac12|1-e^{2v(\delta)}|
\]
and
\[
|v(\delta)|
\ge
\frac{|\delta|}{2}\log\frac{t}{2\pi}.
\]
\end{proof}

This theorem exposes the mechanism visible in the eccentric lobe:
the central circle touches the origin, but transverse continuation changes
the modulus of the rotating component relative to the displaced center.
Exact contact is destroyed immediately for every nonzero displacement.

\begin{figure}[H]
\centering
\includegraphics[width=.70\textwidth]{../figures/fig04_circular_lobe_first_order_no_exit.png}
\caption{Constant-phase-speed version of the same mechanism: the eccentric
center and rotating radius cease to satisfy the tangency equality as soon as
$\delta\ne0$.}
\end{figure}

\begin{figure}[H]
\centering
\includegraphics[width=.82\textwidth]{../figures/fig15_rvm_no_exit_barrier.png}
\caption{Explicit positive RVM no-exit barrier from
Theorem~\ref{thm:rvmnoexit}.}
\end{figure}

This is the exact model result that makes the actual zeta lobe geometry
important: the ideal geometry produced by the same RVM phase already has the
desired first-order confinement mechanism.

\section{Quantifying the eccentricity of actual zeta lobes}

For a rotated exact lobe in uniform phase coordinate define
\[
\boxed{
\mathcal C_n
=
\int_0^1\widehat W_n(q)\,\dd q
}
\]
and
\[
\boxed{
E_n=\frac{2|\mathcal C_n|}{A_n}.
}
\]
An exact tangent circle has $E_n=1$.

Define also
\[
\boxed{
D_n=
\frac1{A_n}
\left[
\int_0^1
\left(
|\widehat W_n(q)-\mathcal C_n|-|\mathcal C_n|
\right)^2
\dd q
\right]^{1/2}.
}
\]
An exact circle centered at $\mathcal C_n$ and tangent to the origin has
$D_n=0$.

For the one-arch model,
\[
\frac{\mathcal C_n^{(\mathrm{arch})}}{A_n}
=
\frac{1+e^{-i\pi s_n}}{\pi(1-s_n^2)}
\]
with continuous value $-i/2$ at $s_n=1$, so
\[
\boxed{
E_n^{(\mathrm{arch})}
=
\frac{4|\cos(\pi s_n/2)|}
{\pi|1-s_n^2|},
\qquad
E_n^{(\mathrm{arch})}\to1\quad(s_n\to1).
}
\]

Thus the phase sweep alone predicts the dominant displacement of the lobe
center.

For the first 49 intervals,
\[
\operatorname{mean}(E_n)=0.989545,
\qquad
\operatorname{median}(E_n)=1.012541,
\]
\[
\operatorname{mean}(D_n)=0.139212,
\qquad
\operatorname{median}(D_n)=0.134423.
\]
The observed organization is strong:
\[
\boxed{\operatorname{corr}(E_n,s_n)=-0.998661,}
\]
\[
\boxed{\operatorname{corr}(D_n,|s_n-1|)=0.989876.}
\]
Most significantly,
\[
\boxed{
\operatorname{corr}(E_n,E_n^{(\mathrm{arch})})=0.997589,
}
\]
with
\[
\boxed{
\operatorname{RMSE}=0.017527,
\qquad
\operatorname{MAE}=0.014155.
}
\]

\begin{figure}[H]
\centering
\includegraphics[width=.84\textwidth]{../figures/fig05_eccentricity_En.png}
\caption{Normalized eccentricity $E_n$ of the first 49 actual lobes.}
\end{figure}

\begin{figure}[H]
\centering
\includegraphics[width=.84\textwidth]{../figures/fig06_circle_tangency_defect_Dn.png}
\caption{Circle-tangency defect $D_n$.}
\end{figure}

\begin{figure}[H]
\centering
\includegraphics[width=.75\textwidth]{../figures/fig19_eccentricity_phase_law.png}
\caption{The measured eccentricity follows the closed-form phase-sweep law
very closely.}
\end{figure}

\begin{figure}[H]
\centering
\includegraphics[width=.72\textwidth]{../figures/fig14_eccentricity_model_vs_actual.png}
\caption{Measured versus one-arch predicted eccentricity.}
\end{figure}

The result is stronger than the statement that the lobes ``look circular''.
Their eccentric center displacement is quantitatively controlled by the same
phase sweep that RVM fixes to one unit in the ideal skeleton.

\section{Euclidean arc geometry}

Define
\[
\boxed{
L_n
=
\int_{\gamma_n}^{\gamma_{n+1}}
\sqrt{
1+
\left|
\zeta'\!\left(\frac12+it\right)
\right|^2
}\,\dd t
}
\]
and
\[
\boxed{
R_n=\frac{L_n}{\gamma_{n+1}-\gamma_n}.
}
\]
Using the Hardy decomposition,
\[
\boxed{
L_n
=
\int_{\gamma_n}^{\gamma_{n+1}}
\sqrt{1+Z'^2+\vartheta'^2Z^2}\,\dd t.
}
\]

This shows why a visually tall lobe can have a much larger length than its
horizontal zero gap: both radial motion and angular sweep contribute.

For example, interval $8\to9$ has
\[
\gamma_8\approx43.32707328,
\qquad
\gamma_9\approx48.00515088,
\]
\[
g_8\approx4.67808,
\qquad
L_8\approx14.508,
\qquad
R_8\approx3.10.
\]
It is long because the thread makes a large transverse excursion while also
sweeping substantial phase.

\begin{figure}[H]
\centering
\includegraphics[width=.87\textwidth]{../figures/fig10_arc_length_Ln.png}
\caption{Directly recomputed Euclidean lobe lengths.}
\end{figure}

\begin{figure}[H]
\centering
\includegraphics[width=.87\textwidth]{../figures/fig11_arc_ratio_Rn.png}
\caption{Arc inflation ratio $R_n=L_n/g_n$.}
\end{figure}

The first-49 correlations are
\[
\operatorname{corr}(R_n,g_n)=0.338036,
\]
\[
\boxed{\operatorname{corr}(R_n,s_n)=0.864608,}
\]
\[
\boxed{\operatorname{corr}(R_n,A_n)=0.911802,}
\]
and
\[
\boxed{\operatorname{corr}(L_n,A_n)=0.995181.}
\]

Thus Euclidean length is overwhelmingly controlled by radial excursion, while
its normalized inflation also retains substantial phase information.

\section{Exact phase-coordinate form of the arc length}

Let
\[
q(t)=\frac{\vartheta(t)}{\pi}.
\]
Since
\[
\dd q=\frac{\vartheta'(t)}{\pi}\dd t,
\]
write $r(q)=|Z(t(q))|$.  Then
\[
\boxed{
L_n
=
\int_{q_n}^{q_{n+1}}
\sqrt{
\left(\frac{\pi}{\vartheta'(t(q))}\right)^2
+
\left(\frac{\dd r}{\dd q}\right)^2
+
\pi^2r(q)^2
}\,\dd q.
}
\]
At the same time,
\[
\boxed{
q_{n+1}-q_n
=
\frac{\vartheta(\gamma_{n+1})-\vartheta(\gamma_n)}{\pi}
=
s_n.
}
\]

Therefore the domain length of the Euclidean lobe integral in phase
coordinates is exactly the phase-unfolded spacing.  This is the direct
geometric bridge between lobe length and zero-spacing statistics.

\section{One-arch elliptic law for arc inflation}

Take
\[
u=\frac{t-\gamma_n}{g_n},
\qquad
Z(t)\approx A_n\sin(\pi u),
\]
and freeze $\vartheta'$ at the midpoint $m_n$.  Set
\[
a_n=\frac{\pi A_n}{g_n},
\qquad
b_n=A_n\vartheta'(m_n).
\]
Then
\[
R_n
\approx
\int_0^1
\sqrt{
1+a_n^2\cos^2(\pi u)+b_n^2\sin^2(\pi u)
}\,\dd u,
\]
hence
\[
\boxed{
R_n^{(\mathrm{arch})}
=
\frac2\pi\sqrt{1+a_n^2}\,
E\!\left(
\frac{a_n^2-b_n^2}{1+a_n^2}
\right),
}
\]
where $E$ is the complete elliptic integral of the second kind.

The recomputed values give
\[
\boxed{
\operatorname{corr}(R_n,R_n^{(\mathrm{arch})})=0.999162,
\qquad
\operatorname{RMSE}=0.116004.
}
\]

\begin{figure}[H]
\centering
\includegraphics[width=.72\textwidth]{../figures/fig12_elliptic_arch_prediction.png}
\caption{Measured arc inflation versus the one-arch elliptic reduction.}
\end{figure}

The same one-arch geometry therefore controls two different observables:
the eccentric displacement of the lobe center and its Euclidean length
inflation.

\section{Angular lobe sweep and local zero statistics}

The angular sweep
\[
\Omega_n
=
\vartheta(\gamma_{n+1})-\vartheta(\gamma_n)
\]
defines
\[
\boxed{s_n=\Omega_n/\pi.}
\]
At large height,
\[
s_n
\sim
\frac{\log(\gamma_n/2\pi)}{2\pi}
(\gamma_{n+1}-\gamma_n).
\]
Thus the conventional local unfolding of zero gaps is geometrically the
angular sweep of consecutive Violet-Thread lobes.

For $k$ consecutive lobes,
\[
\boxed{
\frac{\vartheta(\gamma_{n+k})-\vartheta(\gamma_n)}{\pi}
=
\sum_{j=n}^{n+k-1}s_j.
}
\]
Consequently nearest-neighbor and multi-gap statistics may be read directly
as statistics of one-lobe and cumulative lobe rotations.  No separate
geometric object is required.

\section{Full $2{,}001{,}052$-ordinate audit}

The package contains $2{,}001{,}052$ zeta ordinates and the same number of
RVM/Lambert ordinates.  All are used for the global phase/RVM diagnostics.

\begin{center}
\begin{tabular}{lr}
\toprule
Quantity & Value\\
\midrule
Number of ordinates & $2{,}001{,}052$\\
Number of intervals & $2{,}001{,}051$\\
$\gamma_1$ & $14.134725142$\\
$\gamma_N$ & $1{,}132{,}490.658714411$\\
$\tau_1$ & $14.521346953066$\\
$\tau_N$ & $1{,}132{,}490.722587622$\\
Mean $(\gamma_n-\tau_n)$ & $-1.730666\times10^{-7}$\\
Std. $(\gamma_n-\tau_n)$ & $0.188066347$\\
RMS $(\gamma_n-\tau_n)$ & $0.188066300$\\
Max $|\gamma_n-\tau_n|$ & $0.993710851$\\
Mean $s_n$ & $0.999999964$\\
Std. $s_n$ & $0.407523227$\\
Min $s_n$ & $0.005446783$\\
Max $s_n$ & $3.303473427$\\
1\% quantile $s_n$ & $0.223931915$\\
99\% quantile $s_n$ & $2.078183297$\\
\bottomrule
\end{tabular}
\end{center}

\begin{figure}[H]
\centering
\includegraphics[width=.88\textwidth]{../figures/fig07_full_phase_spacing_histogram.png}
\caption{Phase-unfolded lobe sweep for all $2{,}001{,}051$ consecutive
intervals.}
\end{figure}

\begin{figure}[H]
\centering
\includegraphics[width=.90\textwidth]{../figures/fig08_full_zeta_minus_rvm_lambert.png}
\caption{Difference between zeta and RVM/Lambert ordinates across the complete
data set; every 500th point is displayed for legibility.}
\end{figure}

The full data support the interpretation that RVM fixes the mean phase
architecture, while the arithmetic zero sequence supplies the fluctuating
lobe sweep.


\section{Reciprocal lobe geometry and scattering normal form}\label{sec:reciprocal}

The direct lobe geometry becomes more rigid after inversion.  For a rotated
zero-to-zero lobe
\[
\widehat W_n(\phi)=r_n(\phi)e^{-i\phi},
\qquad
0<\phi<\Omega_n,
\]
define
\[
\boxed{
U_n(\phi)=\frac{1}{\widehat W_n(\phi)}
=
\frac{e^{i\phi}}{r_n(\phi)}.
}
\]
The two contacts of the original lobe with the origin become the two
asymptotic ends of an $\infty\!\to\!\infty$ reciprocal branch.

\subsection{RVM as the flat reciprocal state}

For the unit-sweep one-arch model
\[
\widehat W_{\rm RVM}(\phi)
=
A\sin\phi\,e^{-i\phi},
\qquad 0<\phi<\pi,
\]
one has
\[
\boxed{
U_{\rm RVM}(\phi)
=
\frac1A\cot\phi+\frac{i}{A}.
}
\]
Hence the eccentric tangent circle of the direct picture is exactly a
straight line after inversion.

For a general polar reciprocal curve $\rho(\phi)=1/r(\phi)$, the signed
curvature is
\[
\kappa_U
=
\frac{\rho^2+2\rho_\phi^2-\rho\rho_{\phi\phi}}
     {(\rho^2+\rho_\phi^2)^{3/2}}
=
\frac{r^3(r+r_{\phi\phi})}
     {(r^2+r_\phi^2)^{3/2}}.
\]
Therefore
\[
\boxed{
\operatorname{sgn}\kappa_U
=
\operatorname{sgn}(r_{\phi\phi}+r).
}
\]
In particular,
\[
\boxed{
\kappa_U\equiv0
\iff
r_{\phi\phi}+r=0,
}
\]
and the positive zero-to-zero solution is $r=A\sin\phi$ with total sweep
$\pi$.  Thus the RVM one-arch model is precisely the zero-curvature state of
the reciprocal lobe.

\subsection{Phase spacing as reciprocal scattering angle}

Near the left end,
\[
r_n(\phi)=a_n^-\phi+O(\phi^2),
\]
where
\[
a_n^-=
\frac{|\zeta'(1/2+i\gamma_n)|}{\vartheta'(\gamma_n)}.
\]
Near the right end,
\[
r_n(\phi)=a_n^+(\Omega_n-\phi)+O((\Omega_n-\phi)^2),
\]
with
\[
a_n^+=
\frac{|\zeta'(1/2+i\gamma_{n+1})|}
     {\vartheta'(\gamma_{n+1})}.
\]
Consequently
\[
U_n(\phi)
=
\frac1{a_n^-\phi}
+\frac{i}{a_n^-}+O(\phi)
\]
at the incoming end, whereas at the outgoing end
\[
U_n(\phi)
=
\frac{e^{i\Omega_n}}
     {a_n^+(\Omega_n-\phi)}
-\frac{i e^{i\Omega_n}}{a_n^+}
+O(\Omega_n-\phi).
\]
The scattering-angle defect from the flat RVM state is therefore exactly
\[
\boxed{
\Delta_n^{\rm sc}
=
\Omega_n-\pi
=
\pi(s_n-1).
}
\]

\subsection{Corrected scattering cocycle}

Set
\[
a_n=
\frac{|\zeta'(1/2+i\gamma_n)|}{\vartheta'(\gamma_n)},
\qquad
b_n=\log a_n,
\qquad
\psi_n=\vartheta(\gamma_n)-n\pi.
\]
The endpoint scattering coefficient compatible with the exact lobe
normalization is
\[
\boxed{
\mathfrak S_n
=
\frac{a_{n+1}}{a_n}
e^{\,i(\Omega_n-\pi)}.
}
\]
Since
\[
\Omega_n-\pi=\psi_{n+1}-\psi_n,
\]
the correct complex potential is
\[
\boxed{
Q_n=-b_n-i\psi_n.
}
\]
Then
\[
\boxed{
\mathfrak S_n=e^{Q_n-Q_{n+1}}.
}
\]
Accordingly,
\[
\boxed{
\prod_{j=n}^{n+k-1}\mathfrak S_j
=
\frac{a_{n+k}}{a_n}
\exp\!\left(
i[
\vartheta(\gamma_{n+k})-\vartheta(\gamma_n)-k\pi
]
\right).
}
\]
The real sign in $Q_n$ is essential.  Reversing it would incorrectly replace
$a_{n+1}/a_n$ by its reciprocal.

\subsection{Exact endpoint-normalized lobe factorization}

Using uniform phase coordinate
\[
q=\frac{\phi}{\Omega_n}\in(0,1),
\]
write
\[
\widehat W_n(q)
=
\sin(\pi q)e^{-i\pi s_nq}H_n(q),
\qquad
H_n(q)>0.
\]
The endpoint values are
\[
\boxed{
H_n(0)=s_na_n,
\qquad
H_n(1)=s_na_{n+1}.
}
\]
Define
\[
\boxed{
J_n(q)
=
\frac{H_n(q)}
     {H_n(0)^{1-q}H_n(1)^q}.
}
\]
Then
\[
J_n(0)=J_n(1)=1,
\qquad
J_n(q)>0\quad(0<q<1),
\]
and the exact lobe normal form is
\[
\boxed{
\widehat W_n(q)
=
s_n\sin(\pi q)e^{-i\pi s_nq}
a_n^{1-q}a_{n+1}^{q}J_n(q).
}
\]
All interior information not already carried by the two contacts, the phase
sweep and the endpoint scattering amplitudes is concentrated in the single
residual factor $J_n$.

\subsection{Integrated reciprocal-curvature identity}

Let
\[
\ell_n(q)=\log H_n(q)
\]
on the real lobe.  A direct calculation gives
\[
\boxed{
V_n(q)
:=
\pi^2s_n^2\frac{r_{\phi\phi}+r}{r}
=
\ell_n''
+\ell_n'^2
+2\pi\cot(\pi q)\ell_n'
+\pi^2(s_n^2-1).
}
\]
Multiplying by $\sin^2(\pi q)$ and integrating,
\[
\boxed{
\int_0^1
\sin^2(\pi q)V_n(q)\,\mathrm dq
=
\int_0^1
\sin^2(\pi q)\ell_n'(q)^2\,\mathrm dq
+
\frac{\pi^2}{2}(s_n^2-1).
}
\]
Thus the weighted reciprocal curvature has a strictly positive integral for
every $s_n>1$, independent of the detailed arithmetic deformation.  The
sub-RVM side $s_n<1$ admits competition between the negative RVM term and the
nonnegative deformation energy.

\section{The exponential spectral-purity attempt}\label{sec:spectral-purity}

The preceding normal form suggests an appealing route: prove that the
endpoint-normalized residual $J_n$ is spectrally concentrated so strongly
that it cannot acquire a transverse zero.  This section records the exact
criterion, and also the precise reason why the available analytic input does
not yet prove it.

\subsection{A genuine zero-free criterion}

Suppose, in a complexified affine lobe coordinate $x$, that
\[
J_n(x)-1
=
\sum_{m\ge1} d_{n,m}\sin(m\pi x)
\]
in a band $|\Im x|\le\eta$.  Since
\[
|\sin(m\pi x)|
\le
e^{m\pi\eta},
\]
the condition
\[
\boxed{
\sum_{m\ge1}
|d_{n,m}|e^{m\pi\eta}<1
}
\]
implies
\[
|J_n(x)-1|<1
\]
throughout that band, hence
\[
\boxed{
J_n(x)\neq0.
}
\]
This is a non-circular first-order no-exit criterion: it is a quantitative
property of the same lobe residual, not a restatement of zero-freeness.

\subsection{Why strip analyticity does not prove the sine criterion}

Analyticity of $J_n$ in a complex neighborhood is not enough to force
exponential decay of its ordinary sine coefficients.  Endpoint compatibility
matters.  The entire function
\[
f(x)=x(1-x)
\]
vanishes at both endpoints, but its sine coefficients are
\[
\boxed{
d_m
=
2\int_0^1f(x)\sin(m\pi x)\,\mathrm dx
=
\frac{4(1-(-1)^m)}{\pi^3m^3}.
}
\]
They decay only algebraically.

\begin{figure}[H]
\centering
\includegraphics[width=.78\textwidth]{../figures/fig21_sine_series_endpoint_obstruction.png}
\caption{Even an entire function with zero endpoint values can have only
algebraic sine-series decay.  A strip of analyticity therefore does not by
itself yield the desired weighted sine $\ell^1$ bound.}
\end{figure}

The obstruction is structural: shifting the sine-coefficient contour creates
vertical endpoint contributions unless the odd periodic continuation has the
required compatibility to all orders.  The lobe normalization
$J_n(0)=J_n(1)=1$ does not imply those infinite endpoint identities.

\subsection{A stronger counterexample to real-axis control}

Real-axis closeness, endpoint normalization, and even small ordinary
derivatives do not exclude nearby complex zeros.  For
\[
\boxed{
J_{M,\varepsilon}(x)
=
1+\varepsilon\sin^2(M\pi x)
}
\]
one has
\[
J_{M,\varepsilon}(x)\ge1
\qquad (x\in[0,1]),
\]
and
\[
J_{M,\varepsilon}(0)=J_{M,\varepsilon}(1)=1.
\]
Its complex zeros satisfy
\[
\sin(M\pi x)
=
\pm\frac{i}{\sqrt{\varepsilon}},
\]
hence the nearest zero has distance
\[
\boxed{
|\Im x|
=
\frac{1}{M\pi}
\operatorname{arsinh}\frac1{\sqrt{\varepsilon}}.
}
\]
Taking $\varepsilon=c/M^2$ gives
\[
\sup_{[0,1]}|J_{M,\varepsilon}-1|=O(M^{-2}),
\]
while
\[
|\Im x_{\rm zero}|
\sim
\frac{\log M}{\pi M}
\longrightarrow0.
\]

\begin{figure}[H]
\centering
\includegraphics[width=.78\textwidth]{../figures/fig22_hidden_complex_zeros_counterexample.png}
\caption{A residual can be arbitrarily close to $1$ on the real lobe while
its complex zeros approach the lobe.  Real-axis purity is therefore not the
transverse purity needed for confinement.}
\end{figure}

\subsection{What Phragm\'en--Lindel\"of and convexity can and cannot supply}

A possible analytic input is a uniform bound for the complexified lobe
residual in a fixed neighborhood.  Classical convexity bounds for $\zeta$
can contribute boundary estimates, but four points must be completed before
such an argument is a theorem in the present normalization:

\begin{enumerate}[leftmargin=1.8em]
\item a boundary convexity estimate must be propagated through the whole
strip by a Phragm\'en--Lindel\"of argument with explicit constants;
\item the complex Stirling remainder entering the continuation of
$\vartheta'$ must be bounded uniformly in the required neighborhood;
\item the map from the complex height variable to phase coordinate must be
shown globally univalent on that neighborhood, not merely satisfy
$\vartheta'\neq0$ pointwise;
\item an explicit surviving neighborhood parameter must be computed.
\end{enumerate}

Even if all four items are completed, a bounded analytic strip would still
not establish the preceding sine-series criterion because of the endpoint
obstruction just proved.

\subsection{Chebyshev replacement}

The endpoint obstruction can be removed by using a basis adapted to analytic
functions on a finite interval.  Map $x\in[0,1]$ to $y=2x-1\in[-1,1]$ and
expand
\[
J_n\!\left(\frac{y+1}{2}\right)-1
=
\sum_{m\ge0} c_{n,m}T_m(y).
\]
If the residual extends holomorphically to a Bernstein ellipse
$\mathcal E_\rho$ and
\[
\sup_{\mathcal E_\rho}|J_n-1|\le M,
\]
then the standard contour argument gives
\[
\boxed{
|c_{n,m}|\le 2M\rho^{-m}.
}
\]
Thus a fully explicit strip/ellipse bound would indeed imply exponential
\emph{Chebyshev} spectral decay.

However, exponential decay alone is not enough for no-exit.  On a smaller
ellipse $\mathcal E_r$, $1<r<\rho$, a sufficient condition is
\[
\boxed{
\sum_{m\ge0}|c_{n,m}|\,r^m<1.
}
\]
Convexity and Phragm\'en--Lindel\"of can at best provide an upper bound $M$;
they do not presently provide the required smallness relative to $1$.
Accordingly, the spectral-purity route reaches a genuine analytic
obstruction: the missing statement is a quantitative \emph{closeness} bound
for the normalized residual in a complex neighborhood, not merely
holomorphy or boundedness there.

\subsection{Verdict of the spectral-purity route}

The attempted proof of exponential transverse purity does not close with the
current ingredients.

What is proved:
\begin{enumerate}[leftmargin=1.8em]
\item the endpoint-normalized residual $J_n$ isolates the interior lobe
deformation exactly;
\item a weighted spectral $\ell^1$ bound $<1$ gives a genuine zero-free
band;
\item ordinary sine coefficients need not decay exponentially even for an
entire endpoint-normalized function;
\item an analytic Bernstein-ellipse bound does give exponential Chebyshev
decay;
\item exponential decay without a quantitative smallness bound does not
imply no-exit.
\end{enumerate}

What remains unproved is a non-circular complex-neighborhood estimate strong
enough to force
\[
\sum |c_{n,m}|r^m<1
\]
uniformly for the exact zeta residuals.

Thus the present route ends in a precise failure rather than a hidden
assumption: the real lobe geometry and classical upper bounds do not yet
control the transverse amplitude of the endpoint-normalized residual
strongly enough to exclude complex zeros.


\section{What the geometry has actually established}

The following are exact:
\begin{enumerate}[leftmargin=1.8em]
\item the Euclidean/cylindrical representation of the thread;
\item the zero-to-zero decomposition into lobes;
\item the speed identity
\[
|W_t|^2=Z'^2+\vartheta'^2Z^2;
\]
\item the global first-order equation
\[
F_\delta=-iF_t;
\]
\item the conformal metric and minimal-surface structure;
\item the polar conjugacy
\[
(\log r)_\delta=\alpha_t,
\qquad
\alpha_\delta=-(\log r)_t;
\]
\item the exact tangent-circle identity of the unit-sweep one-arch model;
\item the exact RVM/Lambert tangent-circle skeleton;
\item the exact global no-exit theorem for that RVM model;
\item the phase-coordinate form of the Euclidean arc length;
\item the reciprocal-lobe characterization of RVM as the zero-curvature state;
\item the corrected exact scattering cocycle
\[
\mathfrak S_n=e^{Q_n-Q_{n+1}},
\qquad
Q_n=-\log a_n-i\psi_n;
\]
\item the endpoint-normalized factorization with residual $J_n$;
\item the integrated reciprocal-curvature identity;
\item the rigorous failure of strip analyticity alone to imply exponential
sine spectral purity.
\end{enumerate}

The following are numerical findings from the supplied data:
\begin{enumerate}[leftmargin=1.8em]
\item actual low-height lobes have normalized eccentricity near the tangent
circle value;
\item their eccentricity is predicted almost completely by the phase sweep;
\item the one-arch phase model matches measured $E_n$ with correlation
$0.997589$;
\item the one-arch elliptic model matches measured $R_n$ with correlation
$0.999162$;
\item the global phase-sweep mean is essentially one across all
$2{,}001{,}051$ intervals.
\end{enumerate}

The central analytic target is therefore not an artificial higher-order
expansion.  It is one global statement about the exact first-order geometry:

\[
\boxed{
F(\delta,t)\neq0
\quad
\text{for every nonzero transverse displacement inside the strip}.
}
\]

By Theorem~\ref{thm:firstorder-suff}, a proof of that property from the exact
eccentric lobe geometry closes the zero-confinement problem immediately.

\section{Why the RVM result is the decisive model property}

The RVM model is relevant here for a geometric reason, not merely because
$\tau_n$ approximates $\gamma_n$.

It supplies three properties simultaneously:
\begin{enumerate}[leftmargin=1.8em]
\item exactly one phase unit per lobe;
\item exact eccentric tangent circles;
\item exact global no-exit under the same first-order transverse translation.
\end{enumerate}

The actual zeta lobes then show, numerically, that both their center
displacement and their Euclidean length are extraordinarily well organized
by the corresponding one-arch/RVM phase geometry.

This is the concrete transfer problem generated by the Violet Thread:
identify an exact lobe invariant or inequality, visible in the true
$Z$-deformed lobes, that preserves the RVM no-exit separation under
\[
F_\delta=-iF_t.
\]
Such an invariant would not be a later classical criterion; it would be a
property of the same Euclidean lobe and its exact transverse field.

\section{Provenance of the construction}

The Riemann zeta function, Hardy's $Z$-function, the Riemann--Siegel phase,
the smooth Riemann--von Mangoldt phase, Lambert $W$, Cauchy--Riemann
equations, Euclidean arc length, minimal holomorphic graphs and complete
elliptic integrals are classical ingredients.

The primary construction developed here is the organization of the
three-dimensional value trace into a single geometric object---the Violet
Thread---and the decision to retain its zero-to-zero lobes, Euclidean lengths,
arc inflation ratios, phase sweeps and transverse quarter-turn field as parts
of that same object.  The emphasis on the systematic eccentricity visible in
the actual thread, the phase-centroid observables
$\mathcal C_n,E_n,D_n$, the interpretation of the RVM phase as the generator
of an exact tangent-circle lobe skeleton, and the use of the RVM first-order
no-exit mechanism as the direct geometric model for confinement belong to
this program.

The logical point repeatedly used in the paper is also part of this
construction: once a global no-exit property is obtained from the exact
first-order flow, there is no subsequent hierarchy of higher transverse
orders to check.

\section{Conclusion}

The Violet Thread is not well described as a helix centered on its zero axis.
Its natural pieces are eccentric zero-to-zero lobes.  In uniform phase
coordinate, the basic one-arch geometry becomes a circle tangent to the
origin when the angular sweep is one RVM unit.  The smooth RVM/Lambert
construction produces precisely that tangent-circle architecture.

This observation explains the eccentric location of the zeros in the
Euclidean plot and gives the geometry a dynamical role.  The same exact
first-order equation
\[
F_\delta=-iF_t
\]
that generates the transverse zeta family also generates an exact no-exit
barrier for the RVM tangent circles.  In the actual zeta data, the lobe
eccentricity follows the corresponding RVM-phase prediction with correlation
$0.997589$, while the independent Euclidean arc inflation follows the
one-arch elliptic law with correlation $0.999162$.

The resulting picture is unified:
\[
\boxed{
\text{RVM phase}
\longrightarrow
\text{eccentric tangent-circle skeleton}
\longrightarrow
\text{actual arithmetic lobe deformation}
\longrightarrow
\text{first-order transverse dynamics}.
}
\]
The remaining analytic question is therefore sharply geometric, but the
spectral-purity analysis makes the obstruction more precise.  A bounded
analytic continuation of the endpoint-normalized lobe residual is not enough:
one needs a quantitatively small complex-neighborhood norm, in a basis such
as Chebyshev that is compatible with finite-interval analyticity.  Without
that additional estimate, real-axis circularity, reciprocal curvature,
endpoint scattering and ordinary spectral concentration do not exclude
nearby transverse zeros.  The first-order confinement program therefore
remains open at this explicit norm-control step, not at any hierarchy of
higher transverse orders.

\appendix

\section{Complete first-49 Euclidean arc table}
\small
\begin{longtable}{rrrrrrr}
\toprule
$n$ & $\gamma_n$ & $\gamma_{n+1}$ & $g_n$ & $L_n$ & $R_n$ & $s_n$\\
\midrule
\endfirsthead
\toprule
$n$ & $\gamma_n$ & $\gamma_{n+1}$ & $g_n$ & $L_n$ & $R_n$ & $s_n$\\
\midrule
\endhead
\input{../outputs/arc_metrics_table.tex}
\end{longtable}
\normalsize

\section{Complete first-49 eccentric-lobe table}
\scriptsize
\begin{longtable}{rrrrrrrr}
\toprule
$n$ & $s_n$ & $A_n$ & $\Re\mathcal C_n$ & $\Im\mathcal C_n$ & $E_n$ & $D_n$ & $E_n^{(\mathrm{arch})}$\\
\midrule
\endfirsthead
\toprule
$n$ & $s_n$ & $A_n$ & $\Re\mathcal C_n$ & $\Im\mathcal C_n$ & $E_n$ & $D_n$ & $E_n^{(\mathrm{arch})}$\\
\midrule
\endhead
\input{../outputs/eccentric_lobe_table.tex}
\end{longtable}
\normalsize

\section{Reproducibility package}

The archive contains:
\begin{itemize}[leftmargin=1.6em]
\item \texttt{manuscript/VIOLET\_THREAD\_V3.tex};
\item \texttt{manuscript/VIOLET\_THREAD\_V3.pdf};
\item twenty-two PNG figures;
\item the complete numerical analysis script from Version 2;
\item the Version 3 additional-geometry Python script;
\item the Version 4 analytic-obstruction Python script (which reads no corpus data);
\item \texttt{data/zeross.bin} and \texttt{data/zLamb.bin}, each with
$2{,}001{,}052$ IEEE-754 doubles;
\item CSV outputs for the global audit, exact first-49 arc metrics and
eccentric-lobe metrics;
\item generated TeX tables used in the appendices.
\end{itemize}

\begin{thebibliography}{9}
\bibitem{Titchmarsh}
E. C. Titchmarsh, revised by D. R. Heath-Brown,
\emph{The Theory of the Riemann Zeta-Function}, 2nd ed.,
Oxford University Press, 1986.

\bibitem{Edwards}
H. M. Edwards,
\emph{Riemann's Zeta Function},
Dover Publications, 2001.

\bibitem{Montgomery}
H. L. Montgomery,
The pair correlation of zeros of the zeta function,
\emph{Proc. Sympos. Pure Math.} \textbf{24} (1973), 181--193.

\bibitem{Corless}
R. M. Corless, G. H. Gonnet, D. E. G. Hare, D. J. Jeffrey and D. E. Knuth,
On the Lambert $W$ function,
\emph{Advances in Computational Mathematics} \textbf{5} (1996), 329--359.
\end{thebibliography}

\end{document}
